%
% teil1.tex -- Beispiel-File für das Paper
%
% (c) 2020 Prof Dr Andreas Müller, Hochschule Rapperswil
%
% !TEX root = ../../buch.tex
% !TEX encoding = UTF-8
%
\section{Automatic Differentiation
\label{step:section:teil1}}
\kopfrechts{Automatic Differentiation}
\index{Automatic Differnetiation}%

Aufgrund dem bei der Finite-Differenzen-Methode vorgenommenen Abbruch der Taylor-Reihe entstehen --- wie im vorherigen Abschnitt gezeigt --- Ungenauigkeiten, insbesondere wenn die Schrittweite $h$ zu gross gewählt wird.  
Im Folgenden wird eine Methode vorgestellt, mit der diese Ungenauigkeiten vermieden werden können. Ausgangspunkt ist (\ref{step:taylor2}), in der die Terme höherer Ordnung weiterhin vollständig enthalten sind
\begin{equation*}
	f(x+h) = \sum_{k=0}^{\infty} \frac{h^k}{k!} f^{(k)}(x).
\end{equation*}
Für eine praktisch berechenbare Lösung müssen auch hier die Terme höherer Ordnung eliminiert werden. Im Gegensatz zur Finite-Differenzen-Methode werden diese jedoch nicht einfach abgeschnitten, sondern gezielt zum Verschwinden gebracht.  
Dies wird erreicht, indem für die Schrittweite $h$ spezielle algebraische Eigenschaften eingeführt werden. Anstelle von $h$ wird dabei $\epsilon$ verwendet, um eine klare Abgrenzung zur Finite-Differenzen-Methode zu schaffen.

Sei dazu $\epsilon \neq 0$ ein \emph{nilpotentes} Element, d.h ein Element, für das $\epsilon^2 = 0$ postuliert wird \cite{step:baydin2018}. Dann gilt insbesondere
\index{nilpotent}%
\index{epsilon@$\epsilon$}%

\begin{equation*}
	\epsilon^n = 0 \quad \forall\, n \geq 2.
\end{equation*}

\subsection{Herleitung der dualen Zahl und Motivation}

Wird nun in der Taylor-Entwicklung $h$ durch $\epsilon$ ersetzt und die oben definierten Eigenschaften von $\epsilon$ verwendet, ergibt sich
\begin{equation}
	f(x+\epsilon) = \sum_{k=0}^{\infty} \frac{\epsilon^k}{k!} f^{(k)}(x).
		\label{step:taylorepsilon}
\end{equation}
Aufgrund der Eigenschaft $\epsilon^n = 0$ für alle $n \geq 2$ verschwinden sämtliche Terme höherer Ordnung. Es bleiben lediglich die ersten beiden Terme der Reihe übrig:
\begin{equation}
	f(x+\epsilon) = f(x) + \epsilon f'(x).
		\label{step:autodiffdual}
\end{equation}
Dieser Ausdruck wird als sogenannte \emph{duale Zahl} $\left(f\left(x\right), f'\left(x\right)\right)$ interpretiert. Dabei stellt der Realteil $f(x)$ den Funktionswert dar, während der $\epsilon$-Anteil direkt die Ableitung $f'(x)$ enthält \cite{step:baydin2018}. 
\index{duale Zahl}%

Auf den ersten Blick scheint sich der Ausdruck in Gleichung~\eqref{step:taylorepsilon} kaum von der finiten Differenz in \eqref{step:taylor2} zu unterscheiden.
Der entscheidende Unterschied liegt jedoch in der zugrunde liegenden Algebra:  
Bei der Finite-Differenzen-Methode wird $h$ als reelle Zahl gegen $0$ genommen und die Ableitung wird über einen Grenzübergang approximiert.
Dabei wählt man $h$ sehr klein und hofft, dass der Fehler entsprechend gering bleibt. Mit $\epsilon$ hingegen wird ein Konstrukt eingeführt, bei dem alle Terme höherer Ordnung aufgrund der Eigenschaft $\epsilon^2 = 0$ exakt verschwinden.
Da dies ein algebraisches Konstrukt und keine Approximation ist, entsteht kein Approximationsfehler \cite{step:baydin2018}.
Der wesentliche Vorteil dieser Methode liegt somit darin, dass Ableitungen \emph{exakt} und ohne Abbruchfehler berechnet werden können.
Zusätzlich ermöglicht diese Methode die gleichzeitige Berechnung von Funktionswert und Ableitung in einem einzigen Auswertungsschritt.
Dies macht sie besonders attraktiv für numerische Verfahren, da sie sowohl effizient als auch numerisch stabil ist.



\subsection{Algebraische Herleitung der Rechenregeln für duale Zahlen}

Die Rechenregeln für duale Zahlen lassen sich aus ihrer algebraischen Struktur herleiten \cite{step:baydin2018}. 
\begin{definition}[Duale Zahlen]
	\normalfont
	Die Menge der \emph{dualen Zahlen} ist definiert als
\index{D@$\mathbb{D}$}%
\index{duale Zahl}%
	\begin{equation}
		\mathbb{D} := \{\, u + \varepsilon u' \mid u,u' \in \mathbb{R} \,\},
	\end{equation}
	mit
	\begin{equation}
		\varepsilon^2 = 0
	\end{equation}
\end{definition}

Gegeben seien die beiden dualen Zahlen $(u,u') = u + \varepsilon u'$ und $(v,v') = v + \varepsilon v'$.
Die \emph{Addition} zweier dualer Zahlen sei durch
\begin{equation*}
(u + \varepsilon u') + (v + \varepsilon v') 
	= (u+v) + \varepsilon(u'+v')
\end{equation*}
definiert.
Mit Hilfe der Definition der Addition lässt sich auch die \emph{Subtraktion} zweier dualer Zahlen wie folgt
\begin{equation*}
	(u + \varepsilon u') - (v + \varepsilon v') 
	= (u-v) + \varepsilon(u'-v')
\end{equation*}
definieren.

Die \emph{Multiplikation} zweier dualer Zahlen sei definiert durch
\index{duale Zahl!Multiplikation}%
\begin{align*}
	(u + \varepsilon u')(v + \varepsilon v') 
	&= uv + \varepsilon uv' + \varepsilon u'v + \varepsilon^2 u'v'. \\
	\intertext{Mit $\varepsilon^2 = 0$ folgt}
	&= uv + \varepsilon(uv' + u'v).
\end{align*}

Für die \emph{Division} zweier dualer Zahlen muss in einem ersten Schritt die Inverse einer dualen Zahl definiert werden. 
\index{duale Zahl!Division}%
Diese muss existieren und muss ebenfalls ein Element von $	\mathbb{D} $ sein

\begin{equation*}
(v + \varepsilon v')^{-1} \in \mathbb{D}, \quad \text{für}\quad v\neq0.
\end{equation*}
Unter der Annahme, dass die Inverse einer dualen Zahl wieder eine duale Zahl ist, folgt somit
\begin{equation*}
\frac{1}{v + \varepsilon v'} = a + \varepsilon a'
\end{equation*}
und damit
\begin{equation*}
\left(v + \varepsilon v'\right)\left(a + \varepsilon a'\right) = 1.
\end{equation*}
Mit Hilfe der Regeln für die Addition und Multiplikation zweier dualer Zahlen ergibt sich
\begin{equation*}
av + \varepsilon\left(a'v + av'\right) + \varepsilon^2 a' v' = 1.
\end{equation*}
Mit $\varepsilon^2 = 0$ ergibt sich
\begin{align*}
	av &=1 \\
	a'v + av' &= 0
\end{align*}
und daraus folgend 
\begin{align*}
	a &= \frac{1}{v} \\
	a' &= -\frac{av'}{v}.
\end{align*}
Somit gilt für den Kehrwert der dualen Zahl $v +  \varepsilon v'$
\index{duale Zahl!Kehrwert}%
\begin{equation*}
\frac{1}{v + \varepsilon v'} = a + \varepsilon a' = \frac{1}{v} - \varepsilon\frac{v'}{v^2}
\end{equation*}
Mit Hilfe der Herleitung der Inversen einer dualen Zahl gilt damit für die \emph{Division} zweier dualer Zahlen allgemein
\begin{equation*}
	\frac{u + \varepsilon u'}{v + \varepsilon v'} =  \frac{\left(u + \varepsilon u'\right)\left(v - \varepsilon v' \right)}{v^2} = \frac{uv + \varepsilon\left(u'v - uv'\right) - \varepsilon^2u' v'}{v^2} = \frac{uv + \varepsilon\left(u'v - uv'\right)}{v^2}.
\end{equation*}

Zur Bestimmung der \emph{Potenz} einer dualen Zahl wird die Darstellung
\index{duale Zahl!Potenz}%
\begin{equation*}
	(u,u') = u + \varepsilon u'
\end{equation*}
verwendet. Es gilt ausserdem
\begin{equation*}
	(u,u')^n = (u + \varepsilon u')^n.
\end{equation*}
Unter Verwendung des binomischen Lehrsatzes ergibt sich
\index{binomischer Lehrsatz}%
\begin{equation*}
	(u + \varepsilon u')^n 
	= \sum_{k=0}^{n} \binom{n}{k} u^{n-k} (\varepsilon u')^k.
\end{equation*}
Da $\varepsilon^2 = 0$ gilt, verschwinden alle Terme mit $k \geq 2$, sodass nur die ersten beiden Summanden verbleiben:
%\begin{equation*}
%	(u + \varepsilon u')^n 
%	= u^n + \binom{n}{1} u^{n-1} (\varepsilon u').
%\end{equation*}
\begin{equation*}
	(u + \varepsilon u')^n 
	= u^n + \binom{n}{1} u^{n-1} (\varepsilon u')
	= u^n + \frac{n!}{1!(n-1)!} u^{n-1} (\varepsilon u').
\end{equation*}
Dies vereinfacht sich zu
\begin{equation*}
	(u + \varepsilon u')^n 
	= u^n + \frac{n \cdot (n-1)!}{1 \cdot (n-1)!} u^{n-1} (\varepsilon u')
	= u^n + n\, u^{n-1} (\varepsilon u').
\end{equation*}
Somit ergibt die Potenz einer dualen Zahl mit einem natürlichen Exponenten das Ergebnis
\begin{equation*}
	(u,u')^n = (u^n,\; n u^{n-1} u').
\end{equation*}
Verwendet man statt des binomischen Satzes die binomische Reihe, kann man
\index{Reihe!binomische}%
\index{binomische Reihe}%
die Formel auch für beliebige reelle Exponenten erhalten.

Aus diesen Herleitungen ergeben sich nun wohldefinierte Regeln für das Rechnen mit dualen Zahlen.
\begin{satz}[Rechenregeln für duale Zahlen]
	\normalfont
	Für $u,v,u',v' \in \mathbb{R}$ gelten:
	\begin{align}
		\quad\text{\emph{Addition:}}&&
		(u,u') + (v,v') 
		&= (u+v,\; u'+v'), \quad\\
		\quad\text{\emph{Multiplikation:}}&&
		(u,u') \cdot (v,v') 
		&= (uv,\; uv' + u'v), \quad\\
		\quad\text{\emph{Division:}}&&
		\frac{(u,u')}{(v,v')} 
		&= \biggl(\frac{u}{v},\; \frac{u'v - uv'}{v^2}\biggr), \; v \neq 0, \quad\\
		\quad\text{\emph{Potenz:}}&&
		(u,u')^n 
		&= (u^n,\; n u^{n-1} u'), \; n \in \mathbb{R}.\quad
	\end{align}
	\label{step:autodiffregeln}
\end{satz}

Bemerkenswert ist, dass die jeweils zweite Komponente dieser Rechenregeln genau den aus der klassischen Analysis bekannten Ableitungsregeln entspricht.
Sind $u'$ und $v'$ die Ableitungen von $u$ beziehungsweise $v$, so ergeben sich aus den algebraischen Operationen mit dualen Zahlen unmittelbar die Summen-, Produkt-, Quotienten- und Potenzregel.
Die Ableitungsregeln müssen somit nicht zusätzlich auf die zweite Komponente angewendet werden, sondern sind bereits in der Algebra der dualen Zahlen enthalten.
Wird eine Grösse zusammen mit ihrer Ableitung als duale Zahl $(u,u')$ dargestellt, so liefert jede der oben definierten Operationen in ihrer ersten Komponente den zugehörigen Funktionswert und in ihrer zweiten Komponente deren Ableitung.

Die Menge der dualen Zahlen $\mathbb{D}$ bildet zusammen mit den in Satz \ref{step:autodiffregeln} definierten Regeln eine algebraische Struktur mit den folgenden Eigenschaften:
\begin{itemize}
	\item \emph{Assoziativität:} Sowohl die Addition als auch die Multiplikation sind assoziativ.
	\item \emph{Kommutativität:} Addition und Multiplikation sind kommutativ.
	\item \emph{Distributivität:} Die Multiplikation ist distributiv bezüglich der Addition.
\end{itemize}
Diese Eigenschaften folgen unmittelbar aus den entsprechenden Eigenschaften der reellen Zahlen sowie der Relation $\varepsilon^2 = 0$. Insbesondere werden bei der Multiplikation keine Terme höherer Ordnung als $\varepsilon$ berücksichtigt, wodurch die Struktur konsistent bleibt.
Somit bildet $\mathbb{D}$ einen kommutativen Ring mit Eins, wobei das neutrale Element der Multiplikation durch $1 + \varepsilon \cdot 0$ gegeben ist.


\begin{satz}[Algebraische Struktur der dualen Zahlen]
	Die Menge $\mathbb{D}$ bildet mit den definierten Operationen einen kommutativen Ring mit Eins.
\end{satz}
\index{duale Zahl!kommutativer Ring mit Eins}%
Die Ringaxiome folgen aus den entsprechenden Eigenschaften der reellen Zahlen sowie der Relation $\varepsilon^2 = 0$. Insbesondere bleiben Assoziativität, Kommutativität und Distributivität unter dieser Erweiterung erhalten.
\index{assoziativ}%
\index{distributiv}%
\index{kommutativ}%


\subsection{Die Anwendung der dualen Zahlen auf Polynome}
Gegeben sei das Polynom dritten Grades
\begin{equation*}
	f\left(x\right) =  x^3 + 2x² + 4x + 1.
\end{equation*}
Wird das Polynom für eine allgemeine duale Zahl $(u, u')$ evaluiert, so müssen nur die in Satz \ref{step:autodiffregeln} definierten Rechenregeln angewendet werden. Dabei gilt:
\begin{align*}
	f(u, u') &= (u, u')^3 + 2 (u, u')^2 + 4 (u, u') + 1 \\
	&= (u^3, 3 u^2 u') + 2 (u^2, 2 u u') + (4 u, 4 u') + (1, 0) \\
	&= (u^3 + 2 u^2 + 4 u + 1, \; 3 u^2 u' + 4 u u' + 4 u') \\
	&= (f(u), f'(u) \cdot u').
\end{align*}
Setzt man nun $u' = 1$, entspricht dies der Wahl $\frac{dx}{dx} = 1$, also der Ableitung in Richtung der unabhängigen Variablen $x$. Mit dieser Wahl wird der Dualteil der Zahl exakt die Ableitung der Funktion $f$ an der Stelle $u$ darstellen.
Damit erhält man
\begin{equation*}
	f(u, 1) = (f(u), f'(u)),
\end{equation*}
wobei der Realteil $f(u)$ den Funktionswert liefert und der Dualteil $f'(u)$ automatisch die Ableitung enthält.  
\index{Realteil}%
\index{Dualteil}%
Auf diese Weise können Funktionswert und Ableitung gleichzeitig in einem einzigen Schritt bestimmt werden, ohne dass Grenzwerte gebildet oder Approximationen wie bei der Finite-Differenzen-Methode nötig wären. Anschaulich propagiert die duale Zahl auf Grund ihrer algebraischen Struktur die Ableitung „mit“: Jede Operation auf dem Polynom wendet automatisch die entsprechenden Ableitungsregeln an und liefert den korrekten Ableitungswert im Dualteil, ohne dass ein Grenzübergang oder Approximation wie bei der Finite-Differenzen-Methode erforderlich ist \cite{step:baydin2018}.

Am Beispiel des Polynoms wurde gezeigt, dass sich mit Hilfe dualer Zahlen dessen Ableitung exakt berechnen lässt.
Es treten keine Abbruchfehler auf, da keine Reihen abgeschnitten werden müssen. 
Ebenso entstehen keine Rundungsfehler im Sinne numerischer Differenzenverfahren, da die Ableitung nicht approximiert, sondern exakt über die algebraische Struktur der dualen Zahlen bestimmt wird.
Allerdings ergibt sich eine wichtige Erweiterung: Die bisherigen Rechenregeln lassen sich nicht nur auf polynomiale bzw. algebraische Ausdrücke anwenden, sondern auf alle (reell) differenzierbaren Funktionen, die lokal durch eine Taylor-Reihe beschrieben werden können.
Insbesondere gilt dies für alle analytischen Funktionen \cite{step:baydin2018}.
Damit lassen sich auch elementare Funktionen wie die Exponentialfunktion oder trigonometrische Funktionen problemlos im Rahmen der dualen Zahlen behandeln.
Die Ableitungsregeln folgen dabei weiterhin direkt aus der algebraischen Struktur der dualen Zahlen und der Kettenregel, ohne dass zusätzliche Approximationen notwendig sind.
Im reellen Fall existieren zwar differenzierbare Funktionen, die nicht analytisch sind; im Kontext der dualen Zahlen werden jedoch genau diejenigen Funktionen korrekt erfasst, die eine lokale Taylor-Entwicklung besitzen.
Für komplex differenzierbare Funktionen gilt sogar, dass sie automatisch analytisch sind.
Daher treten in der komplexen Analysis keine nicht-analytischen differenzierbaren Funktionen auf.

\subsection{Anwendung der dualen Zahlen in Python und Vergleich mit der Finite-Differenzen-Methode}
\index{Python}%

Wie bereits gezeigt, lassen sich Polynome mit Hilfe dualer Zahlen exakt ableiten, was im Gegensatz zur Finite-Differenzen-Methode steht, bei der aufgrund der Approximation kleine Abweichungen auftreten.
Im Folgenden wird zusätzlich demonstriert, wie sich auch nicht-algebraische Funktionen mithilfe dualer Zahlen auf ‚automatische‘ Weise ableiten lassen.
Zur Veranschaulichung ist es sinnvoll, eine mögliche Implementierung der dualen Zahlen in Python zu betrachten.
Die vollständige Implementierung ist unter
%\url{https://github.com/smalacar/seminar_algebra_step_code}
\cite{step:pythonstep}
zu finden.
Im Folgenden werden lediglich ausgewählte Ausschnitte gezeigt, gefolgt von Beispielen zur Ableitung sowohl von Polynomen als auch von nicht-algebraischen Funktionen.
In der Klasse \texttt{Dual} werden alle in Satz\,\ref{step:autodiffregeln} dargestellten Regeln mit Hilfe der magischen Methoden implementiert, bspw. die Addition, Multiplikation zweier dualer Zahlen sowie auch die Potenz einer dualen Zahl mit einem reellen Exponenten:
\index{Python!magische Methode}%

\begin{lstlisting}[
	language=Python,
	numbers=none,
	basicstyle=\ttfamily\fontsize{7}{9}\selectfont,
	keepspaces=true,
	]
class Dual:
   def __init__(self, real, dual=0.0):
      self.real = real
      self.dual = dual
	
   def __add__(self, other):
      if isinstance(other, Dual):
         return Dual(self.real + other.real, self.dual + other.dual)
      else:
         return Dual(self.real + other, self.dual)
	
   def __mul__(self, other):
      if isinstance(other, Dual):
         return Dual(
            self.real * other.real,
            self.real * other.dual + self.dual * other.real
         )
      else:
         return Dual(self.real * other, self.dual * other)
	
   def __pow__(self, power):
      return Dual(
         self.real**power,
         power * (self.real ** (power - 1)) * self.dual
      )

\end{lstlisting}

\noindent Mit dieser einfachen Implementierung lassen sich bereits Polynome der Form
\begin{lstlisting}[
	language=Python,
	numbers=none,
	basicstyle=\ttfamily\fontsize{7}{9}\selectfont,
	keepspaces=true,
	]
def f(x):
   return x**3 + 2*x**2 + 4*x + 1
\end{lstlisting}
problemlos ableiten und zwar wie folgt:
\begin{lstlisting}[
	language=Python,
	numbers=none,
	basicstyle=\ttfamily\fontsize{7}{9}\selectfont,
	keepspaces=true,
	]
>>> x = Dual(1, 1)
>>> result = f(x)
>>> result.real
8
>>> result.dual
11
\end{lstlisting}
Der Leser ist eingeladen, die dargestellten Ergebnisse selbständig zu überprüfen.

Im Vergleich dazu wird dieselbe Funktion \texttt{f(x)} mit Hilfe der Methode der finiten Differenzen abgeleitet:
\begin{lstlisting}[
	language=Python,
	numbers=none,
	basicstyle=\ttfamily\fontsize{7}{9}\selectfont,
	keepspaces=true,
	]
def forward_diff(f, x, h):
   return (f(x + h) - f(x)) / h

>>> result_diff = forward_diff(f, 1, 1e-12)
11.000977906405751
\end{lstlisting}
Das Ergebnis verdeutlicht, dass die Ableitung mit der Finite-Differenzen-Methode nicht exakt ist. Ursache hierfür ist der Abbruch der Taylor-Reihe, die der Methode zugrunde liegt, wodurch stets eine kleine Abweichung vom exakten Wert entsteht.


Möchte man nun auch nicht-algebraische Funktionen mit Hilfe der dualen Zahlen ableiten, so müssen diese im Sinne der dualen Zahlen definiert werden. In Python lässt sich dies ganz einfach bewerkstelligen, indem man die entsprechende Funktionen definiert, hier gezeigt am Beispield der Sinusfunktion: 
\begin{lstlisting}[
	language=Python,
	numbers=none,
	basicstyle=\ttfamily\fontsize{7}{9}\selectfont,
	keepspaces=true,
	]
import numpy as np

def sin(d):
   return Dual(np.sin(d.real), np.cos(d.real) * d.dual)
  
>>> x = Dual(2*np.pi, 1)
>>> result = sin(x)
>>> result.dual
1.0
>>> result_diff = forward_diff(np.sin, 2*np.pi, 1e-12)
>>> result_diff
1.000088900582341
\end{lstlisting}
Damit ist gezeigt, wie sich nicht-algebraische Funktionen wie der Sinus mithilfe dualer Zahlen und der damit verbundenen automatischen Differentiation ableiten lassen.
In der Implementierung wird der Ableitungsanteil direkt über die Kettenregel propagiert, indem der Dualteil der Eingabe mit der Ableitung der Funktion multipliziert wird.
Das Ergebnis verdeutlicht den Vorteil gegenüber der Finite-Differenzen-Methode: Die Ableitung von $\sin(x)$ an der Stelle $x = 2\pi$ wird exakt als 1.0 berechnet, während die Finite-Differenzen-Methode einen leicht abweichenden Wert liefert, da hier die Taylor-Reihe nur approximativ ausgewertet wird. 

Bisher wurde gezeigt, dass sich mit dualen Zahlen und automatischer Differentiation sowohl Polynome als auch nicht-algebraische Funktionen effizient und exakt ableiten lassen.
Eine weitere präzise Methode zur Berechnung von Ableitungen numerischer Funktionen ist die Complex-Step-Differentiation.
Sie vermeidet die Rundungsfehler, die bei der Finite-Differenzen-Methode auftreten, und liefert sehr genaue erste Ableitungen, ohne dass ein explizites Dualzahl-Framework notwendig ist.
